
\scanpage{60}
\mspage{16}
près dans la catégorie isotopique $S\mathcal TD_{\mathrm{top}}^\pm$, il faut
\deleted{donc} qu'on ait $A'\subset A^+$ i.e. que les $C_a$ $(a\in A')$ soient effectivement marqués.
De plus, quitte à avoir effectué au préalable des gommages de marquage, ou des surmarquages (on peut tjs s'en tirer par surmarquages uniformes d'un côté et d'autre) on supposera $r(a)=r(a')$, si $a'=\sigma(a)$, $a,a'\in A'$.
Ceci dit, l'opération sera bien déterminée dans la catégorie isotopique, pourvu qu'on \deleted{ait} exige que le marquage soit respecté, et que l'on se donne l'\deleted{\unsure{auto}} antiisomorphisme des torseurs sous $\mathbb Z/r\mathbb Z$ $R_a$ et $R_{a'}$ ($r=r(a)=r(a')$) i.e. un élément de
$R_a\mathbin{\wedge}_{\mathbb Z/r\mathbb Z}R_{a'}$.

\deleted{\unsure{NB}\ \unsure{O}}
Au niveau des graphes pondérés, l'opération consiste à rattacher deux à deux \inserted{par leurs bouts} des arêtes libres --- la pondération des sommets n'est pas changée, et la pondération des arêtes (resp. la structure à repères) suit de façon \deleted{\unsure{non}} évidente (resp. non évidente).
\mspage{59}
